

\section{Structured Coordination Strategy Generation}
\label{section 4}
In this section, we abstract a structured representation for coordination strategy design (Section \ref{sec4.1:Common Structure for Coordination Strategy}) to regularize the ambiguity of natural language \textbf{(R1)}. Based on this structure, a three-stage generation method (Section \ref{sec4.2:Three-stage Strategy Generation}) has been designed to automatically generate an initial coordination strategy based on the user's goal \textbf{(R1)}.
\subsection{Structured Representation for Coordination Strategy}
\label{sec4.1:Common Structure for Coordination Strategy}
To maximize the expressiveness of the structure defined for coordination strategy, E1-4 and we collaboratively survey a corpus of 25 LLM-based Multi-agent collaboration papers and 7 high star open source frameworks \cite{hong2023metagpt, li2023camel,wu2023autogen,CrewAI,AutoAgents,chetDev,chen2023agentverse} for Multi-agent coordination. We analyze the common concepts and structures found in the description for coordination strategy in those papers and projects, based on which we establish a common structure for LLM-based Multi-agent coordination strategy. 

We depict the relationship between key concepts involved in the coordination strategy structure as follows:

\begin{itemize}

\item \textbf{Plan Outline:} Provides a blueprint for the overall collaboration, typically breaking the objectives down into a sequence of \textbf{Tasks} to be carried out one after the other.

\item \textbf{Task}: Takes \textbf{Key Objects} as input and output its target \textbf{Key Object}. The task process specifies how \textbf{Agents} collaboratively finished the task.

\item \textbf{Key Object}: Important intermediate objects during collaboration. 

\item \textbf{Agent}: Intelligent entity that can perform \textbf{Action} according to its observation and \textbf{Instruction}. 

\item \textbf{Action}: The smallest unit of agent behavior observable.
 
\item \textbf{Instruction}: Natural language-based specification that tells agent what/how to do.

\end{itemize}

\subsection{Three-stage Strategy Generation}
\label{sec4.2:Three-stage Strategy Generation}
To help users kick off the exploration, we design a three-stage generation method to provide an initial coordination strategy based on the goal of the user, leveraging the coordination capability of LLM. Details of all the prompts used can be found in our project repository.

\subsubsection{Stage1: Plan Outline Generation}
Given a general description of the goal $g$ provided by the user and a set of initial key objects $\mathcal{I} = \{ko_{1},ko_{2}, ..., ko_{n}\}$, the goal of the Plan Outline Generation stage is to draft a plan outline that decomposes the final goal into a sequence of step tasks $\mathcal{P} = \{t_{1},t_{2}, ..., t_{n}\} = \{t_{i}\}|_{i=0}^{n} $. Here each task $t_{i}$ contains the following attributes:
\begin{itemize}

\item \textbf{Step Name:} A clear and concise name summarizing the step.
\item \textbf{Task Content:} Task description of the current step.
\item \textbf{Input Object List:} The input key objects that will be used in the current step.
\item \textbf{Output Object:} The output key object of the current step.

\end{itemize}

To achieve this, we prompt the LLM to serve as an expert plan outline designer to carefully analyze and decompose the goal provided by the user and output the plan outline $\mathcal{P}$:

\begin{equation}
         \mathcal{P} = \text{LLM} ( g, \mathcal{I}, \texttt{prompt}_\texttt{stage1}) 
\end{equation}


\subsubsection{Stage2: Agent Assignment}
After the plan outline is generated, a team of agents $\mathcal{A}_{i} = \{agent_{1}, agent_{2}, ..., agent_{n\}}$ should be assigned for each task $t_{i}$. Those agents are selected from a board of agent candidates $\mathcal{AB} = \{agent_{1}, agent_{2}, ... agent_{m}\} (\mathcal{A}_{i} \subseteq \mathcal{AB} )$ provided by user. The agents in agent board $\mathcal{AB}$ can be obtained through role prompting\cite{Expertprompting}, LLM fine-tuning\cite{wei2021finetuned}, retrieval-augmented generation (RAG) \cite{lewis2020retrieval}, or even recruitment from an agent store\cite{NexusGPT,GPT-Store,SuperAGI-Marketplace}. Each agent should have a profile to describe its expertise for the coordinator's reference.

To make suitable agent assignment for each task $t_{i}$, we prompt the LLM to serve as an expert manager that analyzes the aspects of ability needed for $t_{i}$ and read the profile for each candidate agent in agent board $\mathcal{AB}$ to output $\mathcal{A}_{i}$:

\begin{equation}
         \mathcal{A}_{i} = \text{LLM} ( g, t_{i}, \mathcal{AB}, \texttt{prompt}_\texttt{stage2}) 
\end{equation}

\begin{figure*}[t]
    \centering
    \includegraphics[width=1\linewidth]{figs/interface.pdf}
    \caption{
    System interface of \textit{AgentCoord}. The first three views (\textit{Plan Outline View}, \textit{Agent Assignment View}, \textit{Task Process View}) correspond to the three-stage coordination strategy generation process while the last view presents the execution result.}
    \label{fig:systemInterface}
\end{figure*}

\subsubsection{Stage3: Task Process Generation}
\label{section 4.2.3}
Once the team of agents $\mathcal{A}_{i}$ are assigned to task $t_{i}$, we can finally specify the task process $\mathcal{S}_{i} = \{action_{1}, action_{2}, ..., action_{n}\}$ for task $t_{i}$, which describes how agents conduct actions to collaboratively finish task $t_{i}$. Here each action contains the following attributes:

\begin{itemize}

\item \textbf{Agent Name:} The name of the agent to conduct this action.
\item \textbf{Instruction:} The instruction for this action, tells the agent what/how to do.
\item \textbf{Interaction Type:} Classify the action based on cooperative interaction type, which can be one of  ``propose'' (propose something that may contribute to the current task), ``critique'' (provide feedback to the action result of other agents), ``improve'' (improve the result of a previous action), and ``finalize'' (deliver the final result for current task based on previous actions).
\item \textbf{Important Input:} Previous information that is important for performing current action, which can be certain action results of other agents or previous key objects.

\end{itemize}

Note that although it is enough to use ``Agent Name'' and ``Description'' to define an action for execution, here we propose two auxiliary attributes (``Interaction Type'' and ``Important Input'') to help articulate its relationship with other actions in the context of collaboration. 

To generate the specification for the task process $t_{i}$, we prompt the LLM to serve as an expert collaboration coordinator to carefully read the profile of each agent assigned to the current task $t_{i}$ and output the task process specification ${S}_{i}$:


\begin{equation}
         \mathcal{S}_{i} = \text{LLM} ( g, t_{i}, \mathcal{A}_{i}, \texttt{prompt}_\texttt{stage3}) 
\end{equation}

